\section{\texorpdfstring{\(L^2\) 的几何让投影与正交可用}{L2 的几何让投影与正交可用}}
\label{sec:12-Real-Analysis-Support/07-Lp-Spaces-and-Function-Geometry/03}

\subsection{困境：逼近一个函数，什么叫``最好''}

手里有一个函数 \(f\)，想用一族简单函数（多项式、三角函数、台阶函数……）的线性组合来逼近它。有限维线性空间 \(M\) 已经选好——比如次数不超过 \(N\) 的多项式——问题是：怎么挑组合系数，让逼近``最好''？

在 \(\R^n\) 里人人都会：做最小二乘，残差垂直于逼近子空间。背后的几何画面干干净净：投影算子、正交分解、Pythagoras 定理。可函数空间不像 \(\R^n\)——函数之间的距离有不止一种量法。你可以量最大逐点偏差（\(L^\infty\)），量绝对偏差的积分（\(L^1\)），也可以量平方偏差的积分（\(L^2\)）。不同量法给出不同的``最近点''，而且大部分量法给不出正交投影那种漂亮的线性结构。

\subsection{试着照搬有限维的配方}

回忆 \(\R^n\) 里最小二乘为什么好用。你有一个内积 \(\langle x,y\rangle=x^{\trans}y\)，范数是内积自己诱导的 \(\norm{x}=\sqrt{\langle x,x\rangle}\)。给一个子空间 \(M\)，最近点 \(Px\) 存在且唯一，而且误差 \(x-Px\) 垂直于 \(M\) 中任何向量。Pythagoras 定理说 \(\norm{x}^2=\norm{Px}^2+\norm{x-Px}^2\)——总能量的正交分解。

要把这套搬到函数空间，函数之间的``内积''长什么样？最自然的候选是

\[
\langle f,g\rangle=\int f(x)g(x)\,d\mu(x).
\]

它对称、双线性、正定（\(\langle f,f\rangle=0\) 推出 \(f=0\) 几乎处处）。配这个内积的范数是

\[
\norm{f}_2=\sqrt{\langle f,f\rangle}=\left(\int |f|^2\,d\mu\right)^{1/2},
\]

正是 \(L^2\) 范数。在这个范数下把 \(L^2\) 中所有 Cauchy 列收拢完备化，就得到了 \(L^2\) 空间——一个 Hilbert 空间。

\subsection{碰壁：为什么不是 \(L^1\) 或 \(L^\infty\)}

\(L^1\) 的范数 \(\norm{f}_1=\int|f|\,d\mu\) 不能由任何内积诱导。\(L^\infty\) 的范数 \(\norm{f}_\infty=\operatorname{ess\,sup}|f|\) 也不行。判别法则是平行四边形法则：

\[
\norm{f+g}^2+\norm{f-g}^2=2(\norm{f}^2+\norm{g}^2).
\]

一个范数能由内积诱导，当且仅当它满足这条法则。在所有 \(L^p\) 中，只有 \(p=2\) 通过这项测试。

后果是根本性的：在 \(L^1\) 或 \(L^\infty\) 中，``角度''和``正交''这两个词没有定义。你可以谈距离、谈收敛、谈最佳逼近的存在性（\(L^p\) 一致凸保证唯一性），但你不能说``误差垂直于逼近子空间''——因为没有垂直这个概念。逼近算子一般是非线性的，没有 Pythagoras 分解，没有 Parseval 等式。Fourier 分析、最小二乘和量子力学都天然栖息在 \(L^2\) 里，原因不是 \(L^2\) 范数更``自然''，而是它是唯一自备内积的 \(L^p\)。

\subsection{精确构造一：正交投影定理}

设 \(M\subseteq L^2\) 是闭线性子空间。对任意 \(f\in L^2\)，在 \(M\) 中存在唯一的点 \(Pf\) 使

\[
\norm{f-Pf}_2=\min_{m\in M}\norm{f-m}_2,
\]

并且误差与整个子空间正交：

\[
\langle f-Pf,m\rangle=0\quad\text{对所有 }m\in M.
\]

证明与有限维最小二乘的推导同构。若误差在 \(M\) 中还有非零分量，沿该分量方向微调 \(Pf\) 就能继续缩短到 \(f\) 的距离——完备性保证你能在 \(M\) 中取到极限。闭性不可省略：逼近序列的极限必须仍然属于 \(M\)，否则``最近点''可能溜到子空间外面去。有限维子空间自动闭；无限维子空间要单独验证闭性。

Pythagoras 定理随之而来：

\[
\norm{f}_2^2=\norm{Pf}_2^2+\norm{f-Pf}_2^2.
\]

总能量 = 子空间内的部分 + 正交残差。一切 \(L^2\) 中的投影分解——Fourier 级数、小波展开、条件期望——都是这个恒等式的变体。

\subsection{精确构造二：Riesz 表示定理}

设 \(\ell\) 是 \(L^2\) 上的连续线性泛函（一个把函数映到标量的线性映射，且 \(\ell(f_n)\to\ell(f)\) 只要 \(f_n\to f\) 在 \(L^2\) 中）。则存在唯一的 \(u\in L^2\)，使

\[
\ell(f)=\langle f,u\rangle\quad\text{对所有 }f\in L^2.
\]

有限维是常识：每个线性泛函由一组系数决定，把那组系数排成向量，泛函就是与那个向量的内积。Hilbert 空间的证明用正交分解走同一条路：若 \(\ell\neq0\)，其核 \(\ker\ell\) 是真闭子空间，正交补是 1 维的，从补中取一个向量适当缩放即得表示元 \(u\)。连续性的作用是保证核是闭的——没有连续性，正交补可能退化，\(u\) 就找不出来。

Riesz 表示定理直接解释了 RKHS（再生核 Hilbert 空间）为什么特殊。求值泛函 \(\delta_x:f\mapsto f(x)\) 在一般的 \(L^2\) 中不连续——你可以改变函数在一个零测集上的值而 \(L^2\) 范数纹丝不动，但逐点求值却剧烈跳变。所以在 \(L^2\) 中没有与 \(\delta_x\) 对应的 Riesz 表示元。RKHS 的定义正是：要求所有求值泛函连续，从而每个 \(x\) 有对应的 Riesz 表示元 \(k(x,\cdot)\)，满足再生性质 \(\langle f,k(x,\cdot)\rangle_{\mathcal H_k}=f(x)\)。这一性质不是核方法的设计选择——它是``要求求值泛函连续''这个几何条件在 Hilbert 空间框架下的必然推论。\(L^2\) 和 RKHS 的全部差别，浓缩在这一句话里。

\subsection{数值落地：Fourier 级数是 \(L^2\) 的正交展开}

在 \(L^2([-\pi,\pi])\) 上，复指数函数系

\[
e_k(x)=\frac{1}{\sqrt{2\pi}}e^{ikx},\qquad k\in\Z
\]

是标准正交的：\(\langle e_k,e_\ell\rangle=\delta_{k\ell}\)。取前 \(2N+1\) 个基函数张成的有限维子空间 \(M_N\)，把 \(f\) 向 \(M_N\) 投影。由投影定理，系数恰好是内积

\[
\widehat f(k)=\langle f,e_k\rangle=\frac{1}{\sqrt{2\pi}}\int_{-\pi}^{\pi}f(x)e^{-ikx}\,dx,
\]

部分和

\[
S_Nf=\sum_{|k|\le N}\widehat f(k)e_k
\]

是 \(M_N\) 中对 \(f\) 的最近点——在所有不超过 \(N\) 次的正弦波的线性组合中，Fourier 部分和在 \(L^2\) 意义下是最优的。

Bessel 不等式是 Pythagoras 定理的直接推论：

\[
\sum_{|k|\le N}|\widehat f(k)|^2\le\norm{f}_2^2.
\]

有限截断的能量永远不会超过总能量。三角函数系在 \(L^2([-\pi,\pi])\) 中完备——它们张成的子空间的并集稠密——因此 \(S_Nf\to f\) 在 \(L^2\) 中，Bessel 升级为 Parseval 恒等式：

\[
\sum_{k\in\Z}|\widehat f(k)|^2=\norm{f}_2^2.
\]

``时域总能量等于频域总能量''不是额外的物理定律，是标准正交基完备时 Pythagoras 定理的自然结果。把复指数基换成 Legendre 多项式、Haar 小波或任何完备标准正交系，恒等式一个字不改。

\subsection{这套几何只在 \(L^2\) 中免费}

其他 \(L^p\) 没有内积，正交失去定义。最佳逼近仍存在（\(L^p\) 对 \(1<p<\infty\) 是一致凸的，保证唯一性），但逼近算子一般不是线性的，误差没有``垂直于子空间''的干净刻画，级数展开没有 Bessel 和 Parseval 这样的等式。

Fourier 分析、量子力学的态空间、最小二乘回归和条件期望全部住在 \(L^2\) 中，这不是巧合——是只有那里有一整套正交几何工具免费用。

\subsection{落点}

\(L^2\) 的三件几何工具——投影定理、Riesz 表示、正交展开——全部由``内积 + 完备''两个条件驱动。下一节把它们用到概率论最核心的对象上：条件期望。它已由局部积分匹配定义；在 \(L^2\) 中，它获得了更锐利的身份——向 \(\sigma\)-代数可测函数子空间的正交投影。
